\section{十九世纪战争中的省份、家庭与物质生活}
太平、捻军、西北战争和条约战争改变了省级财政、河流、城市、宗教组织和家庭迁徙。本节拒绝把战报的伤亡、地方官的群体标签或战后企业的设立当成社会结果本身，并将地方档、契约、家谱、报刊、教会与外文记录作为必要的交叉约束。
\subsection{1850—1859：地方材料所能看见的尺度}
以下事件按账本时间排列；它们不是同质的社会样本，而是提示应到何种地点、机构和材料中继续追问的入口。
\hypertarget{social:EQ-1851-JINTIAN-UPRISING}{}\paragraph{咸丰元年（1851年）}洪秀全集团在金田起事，太平天国战争开始。从地方治理、身份、家庭与物质生活的角度，这一节点要求追问它在何处改变了资源、道路、组织或日常风险。核心来源为《清文宗实录》咸丰元年相关条；军机处档、战事方略及地方奏报。这类编年材料可以标定朝廷获知和叙述事件的时间，却往往压低妇女、儿童、佃户、流民和非文书精英的行动。账本所示的研究方向是战争动员、军需、战场暴力与地方社会研究。可核的条件包括宗教组织、地方社会冲突和清末财政治理危机交织。应分层观察的后果是参与者动机与人数不能由官府或太平军单方叙事概括。证据界限：以咸丰元年（1851年）的同期文书为定位起点；官修记录、奏报或外交文本服务特定机构立场，具体日期、规模、地方执行与对手视角须以独立材料复核。
\hypertarget{social:EQ-1851-YONGAN-KINGSHIPS}{}\paragraph{咸丰元年（1851年）}太平军在永安建立王号和军事组织。从地方治理、身份、家庭与物质生活的角度，这一节点要求追问它在何处改变了资源、道路、组织或日常风险。核心来源为《清文宗实录》咸丰元年相关条；宫中朱批奏折、内阁大库题本与《清史稿》相关志传。它首先留下的是官府分类、任命、处置或资源调拨的痕迹；要判断基层是否照办，仍须寻找县档、账簿、契约或后续案件。账本所示的研究方向是宫廷政治、制度运作与中央—地方信息链研究。可核的条件包括起事扩展需要更稳定的指挥与合法性语言。应分层观察的后果是太平天国制度在流动战时环境中不断变化。证据界限：以咸丰元年（1851年）的同期文书为定位起点；官修记录、奏报或外交文本服务特定机构立场，具体日期、规模、地方执行与对手视角须以独立材料复核。
\hypertarget{social:EQ-1852-TAIPING-NORTHWARD-MARCH}{}\paragraph{咸丰二年（1852年）}太平军北上经湖南，清军防御与地方团练开始加速组织。从地方治理、身份、家庭与物质生活的角度，这一节点要求追问它在何处改变了资源、道路、组织或日常风险。核心来源为《清文宗实录》咸丰二年相关条；军机处档、战事方略及地方奏报。这类编年材料可以标定朝廷获知和叙述事件的时间，却往往压低妇女、儿童、佃户、流民和非文书精英的行动。账本所示的研究方向是战争动员、军需、战场暴力与地方社会研究。可核的条件包括广西战场难以封堵太平军的机动路线。应分层观察的后果是地方武装的出现改变后续战争权力结构。证据界限：以咸丰二年（1852年）的同期文书为定位起点；官修记录、奏报或外交文本服务特定机构立场，具体日期、规模、地方执行与对手视角须以独立材料复核。
\hypertarget{social:EQ-1852-HUNAN-ARMY-ORIGINS}{}\paragraph{咸丰二年（1852年）}曾国藩等在湖南组织团练和湘军雏形。从地方治理、身份、家庭与物质生活的角度，这一节点要求追问它在何处改变了资源、道路、组织或日常风险。核心来源为《清文宗实录》咸丰二年相关条；军机处档、战事方略及地方奏报。这类编年材料可以标定朝廷获知和叙述事件的时间，却往往压低妇女、儿童、佃户、流民和非文书精英的行动。账本所示的研究方向是战争动员、军需、战场暴力与地方社会研究。可核的条件包括八旗、绿营难以独力应对内战。应分层观察的后果是地方募兵和筹饷为清廷提供新力量，也提高地方军政自主性。证据界限：以咸丰二年（1852年）的同期文书为定位起点；官修记录、奏报或外交文本服务特定机构立场，具体日期、规模、地方执行与对手视角须以独立材料复核。
\hypertarget{social:EQ-1853-NANJING-TAKING}{}\paragraph{咸丰三年（1853年）}太平军攻占南京并改名天京，建立长期政权中心。从地方治理、身份、家庭与物质生活的角度，这一节点要求追问它在何处改变了资源、道路、组织或日常风险。核心来源为《清文宗实录》咸丰三年相关条；军机处档、战事方略及地方奏报。这类编年材料可以标定朝廷获知和叙述事件的时间，却往往压低妇女、儿童、佃户、流民和非文书精英的行动。账本所示的研究方向是战争动员、军需、战场暴力与地方社会研究。可核的条件包括长江交通和清军防线崩解提供战略机会。应分层观察的后果是江南财政文化中心易手，战争波及全国资源网络。证据界限：以咸丰三年（1853年）的同期文书为定位起点；官修记录、奏报或外交文本服务特定机构立场，具体日期、规模、地方执行与对手视角须以独立材料复核。
\hypertarget{social:EQ-1853-NORTHERN-EXPEDITION}{}\paragraph{咸丰三年（1853年）}太平军发动北伐，进逼华北但补给与指挥困难加重。从地方治理、身份、家庭与物质生活的角度，这一节点要求追问它在何处改变了资源、道路、组织或日常风险。核心来源为《清文宗实录》咸丰三年相关条；军机处档、战事方略及地方奏报。这类编年材料可以标定朝廷获知和叙述事件的时间，却往往压低妇女、儿童、佃户、流民和非文书精英的行动。账本所示的研究方向是战争动员、军需、战场暴力与地方社会研究。可核的条件包括天京建立后希望扩大战略压力。应分层观察的后果是北伐失败限制太平天国对北方的直接威胁。证据界限：以咸丰三年（1853年）的同期文书为定位起点；官修记录、奏报或外交文本服务特定机构立场，具体日期、规模、地方执行与对手视角须以独立材料复核。
\hypertarget{social:EQ-1853-WESTERN-EXPEDITION}{}\paragraph{咸丰三年（1853年）}太平军西征争夺长江中上游，战事牵动安徽、江西、湖北。从地方治理、身份、家庭与物质生活的角度，这一节点要求追问它在何处改变了资源、道路、组织或日常风险。核心来源为《清文宗实录》咸丰三年相关条；军机处档、战事方略及地方奏报。这类编年材料可以标定朝廷获知和叙述事件的时间，却往往压低妇女、儿童、佃户、流民和非文书精英的行动。账本所示的研究方向是战争动员、军需、战场暴力与地方社会研究。可核的条件包括天京需要控制粮食、兵员和长江交通。应分层观察的后果是区域战场长期化，地方社会承受征发与迁徙。证据界限：以咸丰三年（1853年）的同期文书为定位起点；官修记录、奏报或外交文本服务特定机构立场，具体日期、规模、地方执行与对手视角须以独立材料复核。
\hypertarget{social:EQ-1853-SMALL-SWORDS-SHANGHAI}{}\paragraph{咸丰三年（1853年）}上海小刀会起事，租界和华人城市空间关系发生变化。从地方治理、身份、家庭与物质生活的角度，这一节点要求追问它在何处改变了资源、道路、组织或日常风险。核心来源为《清文宗实录》咸丰三年相关条；地方志、督抚奏折及地方档案。这类编年材料可以标定朝廷获知和叙述事件的时间，却往往压低妇女、儿童、佃户、流民和非文书精英的行动。账本所示的研究方向是地方档案、迁徙、宗教、族群与社会动员研究。可核的条件包括城市商业、秘密结社和地方治理矛盾交织。应分层观察的后果是租界保护与地方镇压的边界需结合中外档案分析。证据界限：以咸丰三年（1853年）的同期文书为定位起点；官修记录、奏报或外交文本服务特定机构立场，具体日期、规模、地方执行与对手视角须以独立材料复核。
\hypertarget{social:EQ-1854-TAIPING-NORTHERN-DEFEAT}{}\paragraph{咸丰四年（1854年）}太平军北伐受挫，清军逐步恢复华北防线。从地方治理、身份、家庭与物质生活的角度，这一节点要求追问它在何处改变了资源、道路、组织或日常风险。核心来源为《清文宗实录》咸丰四年相关条；军机处档、战事方略及地方奏报。这类编年材料可以标定朝廷获知和叙述事件的时间，却往往压低妇女、儿童、佃户、流民和非文书精英的行动。账本所示的研究方向是战争动员、军需、战场暴力与地方社会研究。可核的条件包括长距离行军、冬季补给与地方支持不足限制北伐。应分层观察的后果是太平天国转而依赖长江流域与地方战场。证据界限：以咸丰四年（1854年）的同期文书为定位起点；官修记录、奏报或外交文本服务特定机构立场，具体日期、规模、地方执行与对手视角须以独立材料复核。
\hypertarget{social:EQ-1854-XIANG-ARMY-ORGANIZATION}{}\paragraph{咸丰四年（1854年）}湘军编制、饷源和将领网络逐步成形。从地方治理、身份、家庭与物质生活的角度，这一节点要求追问它在何处改变了资源、道路、组织或日常风险。核心来源为《清文宗实录》咸丰四年相关条；宫中朱批奏折、内阁大库题本与《清史稿》相关志传。它首先留下的是官府分类、任命、处置或资源调拨的痕迹；要判断基层是否照办，仍须寻找县档、账簿、契约或后续案件。账本所示的研究方向是宫廷政治、制度运作与中央—地方信息链研究。可核的条件包括清廷需要可用的地方军事力量。应分层观察的后果是军事私人网络与国家授权之间的界线长期模糊。证据界限：以咸丰四年（1854年）的同期文书为定位起点；官修记录、奏报或外交文本服务特定机构立场，具体日期、规模、地方执行与对手视角须以独立材料复核。
\hypertarget{social:EQ-1855-YELLOW-RIVER-COURSE-SHIFT}{}\paragraph{咸丰五年（1855年）}黄河在铜瓦厢决口改道北流，重塑华北水系与社会经济。从地方治理、身份、家庭与物质生活的角度，这一节点要求追问它在何处改变了资源、道路、组织或日常风险。核心来源为《清文宗实录》咸丰五年相关条；地方志、督抚奏折及地方档案。这类编年材料可以标定朝廷获知和叙述事件的时间，却往往压低妇女、儿童、佃户、流民和非文书精英的行动。账本所示的研究方向是地方档案、迁徙、宗教、族群与社会动员研究。可核的条件包括长期河工压力和水文变化共同作用。应分层观察的后果是河道变迁的时间、影响范围和人口后果须依地方资料分区考察。证据界限：以咸丰五年（1855年）的同期文书为定位起点；官修记录、奏报或外交文本服务特定机构立场，具体日期、规模、地方执行与对手视角须以独立材料复核。
\hypertarget{social:EQ-1855-NIAN-REBELLION-EXPANSION}{}\paragraph{咸丰五年（1855年）}捻军在华北、安徽等地活动扩大，形成与太平战争交织的战场。从地方治理、身份、家庭与物质生活的角度，这一节点要求追问它在何处改变了资源、道路、组织或日常风险。核心来源为《清文宗实录》咸丰五年相关条；军机处档、战事方略及地方奏报。这类编年材料可以标定朝廷获知和叙述事件的时间，却往往压低妇女、儿童、佃户、流民和非文书精英的行动。账本所示的研究方向是战争动员、军需、战场暴力与地方社会研究。可核的条件包括灾荒、流民、盐枭和地方武装网络提供条件。应分层观察的后果是捻军组织结构和社会基础不应由“匪”一词简化。证据界限：以咸丰五年（1855年）的同期文书为定位起点；官修记录、奏报或外交文本服务特定机构立场，具体日期、规模、地方执行与对手视角须以独立材料复核。
\hypertarget{social:EQ-1856-TIANJING-INCIDENT}{}\paragraph{咸丰六年（1856年）}天京事变中太平天国领导层发生内讧和大规模杀戮。从地方治理、身份、家庭与物质生活的角度，这一节点要求追问它在何处改变了资源、道路、组织或日常风险。核心来源为《清文宗实录》咸丰六年相关条；宫中朱批奏折、内阁大库题本与《清史稿》相关志传。它首先留下的是官府分类、任命、处置或资源调拨的痕迹；要判断基层是否照办，仍须寻找县档、账簿、契约或后续案件。账本所示的研究方向是宫廷政治、制度运作与中央—地方信息链研究。可核的条件包括权力集中、军政分工与宗教权威冲突累积。应分层观察的后果是事件严重削弱太平天国协调能力，细节和死亡数存在争议。证据界限：以咸丰六年（1856年）的同期文书为定位起点；官修记录、奏报或外交文本服务特定机构立场，具体日期、规模、地方执行与对手视角须以独立材料复核。
\hypertarget{social:EQ-1856-SECOND-OPIUM-WAR-BEGIN}{}\paragraph{咸丰六年（1856年）}亚罗号事件等冲突后英法对清开战，第二次鸦片战争开始。从地方治理、身份、家庭与物质生活的角度，这一节点要求追问它在何处改变了资源、道路、组织或日常风险。核心来源为《清文宗实录》咸丰六年相关条；军机处档、战事方略及地方奏报。这类编年材料可以标定朝廷获知和叙述事件的时间，却往往压低妇女、儿童、佃户、流民和非文书精英的行动。账本所示的研究方向是战争动员、军需、战场暴力与地方社会研究。可核的条件包括条约修订、领事权与商业扩张要求升级。应分层观察的后果是战争与内战并行，迫使清廷分散军事和财政资源。证据界限：以咸丰六年（1856年）的同期文书为定位起点；官修记录、奏报或外交文本服务特定机构立场，具体日期、规模、地方执行与对手视角须以独立材料复核。
\hypertarget{social:EQ-1857-GUANGZHOU-FALL}{}\paragraph{咸丰七年（1857年）}英法联军攻占广州，叶名琛被俘。从地方治理、身份、家庭与物质生活的角度，这一节点要求追问它在何处改变了资源、道路、组织或日常风险。核心来源为《清文宗实录》咸丰七年相关条；军机处档、战事方略及地方奏报。这类编年材料可以标定朝廷获知和叙述事件的时间，却往往压低妇女、儿童、佃户、流民和非文书精英的行动。账本所示的研究方向是战争动员、军需、战场暴力与地方社会研究。可核的条件包括珠江口海军优势和清军防御薄弱。应分层观察的后果是广东地方治理、商贸和对外关系受到持续冲击。证据界限：以咸丰七年（1857年）的同期文书为定位起点；官修记录、奏报或外交文本服务特定机构立场，具体日期、规模、地方执行与对手视角须以独立材料复核。
\hypertarget{social:EQ-1858-TIANJIN-TREATIES}{}\paragraph{咸丰八年（1858年）}《天津条约》等在战争压力下签订，扩大通商、使节与传教权利。从地方治理、身份、家庭与物质生活的角度，这一节点要求追问它在何处改变了资源、道路、组织或日常风险。核心来源为《清文宗实录》咸丰八年相关条；《筹办夷务始末》相关年卷与中外照会、条约文本。它首先留下的是官府分类、任命、处置或资源调拨的痕迹；要判断基层是否照办，仍须寻找县档、账簿、契约或后续案件。账本所示的研究方向是条约文本、外交档案、口岸社会与国际法实践研究。可核的条件包括联军北上威胁京师迫使清廷议约。应分层观察的后果是签署不等于立即批准或实施，后续冲突仍在发生。证据界限：以咸丰八年（1858年）的同期文书为定位起点；官修记录、奏报或外交文本服务特定机构立场，具体日期、规模、地方执行与对手视角须以独立材料复核。
\hypertarget{social:EQ-1859-TAKU-BATTLE}{}\paragraph{咸丰九年（1859年）}大沽口战事使英法换约受阻，战争再度升级。从地方治理、身份、家庭与物质生活的角度，这一节点要求追问它在何处改变了资源、道路、组织或日常风险。核心来源为《清文宗实录》咸丰九年相关条；军机处档、战事方略及地方奏报。这类编年材料可以标定朝廷获知和叙述事件的时间，却往往压低妇女、儿童、佃户、流民和非文书精英的行动。账本所示的研究方向是战争动员、军需、战场暴力与地方社会研究。可核的条件包括条约批准和使节入京问题未获解决。应分层观察的后果是清方短期胜利未改变总体军力与外交劣势。证据界限：以咸丰九年（1859年）的同期文书为定位起点；官修记录、奏报或外交文本服务特定机构立场，具体日期、规模、地方执行与对手视角须以独立材料复核。
\subsection{1860—1869：地方材料所能看见的尺度}
以下事件按账本时间排列；它们不是同质的社会样本，而是提示应到何种地点、机构和材料中继续追问的入口。
\hypertarget{social:EQ-1860-BEIJING-CAMPAIGN}{}\paragraph{咸丰十年（1860年）}英法联军进逼北京，圆明园遭焚掠，咸丰帝出走热河。从地方治理、身份、家庭与物质生活的角度，这一节点要求追问它在何处改变了资源、道路、组织或日常风险。核心来源为《清文宗实录》咸丰十年相关条；军机处档、战事方略及地方奏报。这类编年材料可以标定朝廷获知和叙述事件的时间，却往往压低妇女、儿童、佃户、流民和非文书精英的行动。账本所示的研究方向是战争动员、军需、战场暴力与地方社会研究。可核的条件包括大沽后战争升级与清军防线崩溃。应分层观察的后果是战争记忆、掠夺和城市影响需结合中外见证与物质证据。证据界限：以咸丰十年（1860年）的同期文书为定位起点；官修记录、奏报或外交文本服务特定机构立场，具体日期、规模、地方执行与对手视角须以独立材料复核。
\hypertarget{social:EQ-1860-CONVENTION-OF-PEKING}{}\paragraph{咸丰十年（1860年）}《北京条约》及相关中俄条约确认和扩大列强权益并涉及东北边界。从地方治理、身份、家庭与物质生活的角度，这一节点要求追问它在何处改变了资源、道路、组织或日常风险。核心来源为《清文宗实录》咸丰十年相关条；《筹办夷务始末》相关年卷与中外照会、条约文本。它首先留下的是官府分类、任命、处置或资源调拨的痕迹；要判断基层是否照办，仍须寻找县档、账簿、契约或后续案件。账本所示的研究方向是条约文本、外交档案、口岸社会与国际法实践研究。可核的条件包括北京陷落迫使清廷接受新安排。应分层观察的后果是各条约文本、边界实施和土地控制不可混为一谈。证据界限：以咸丰十年（1860年）的同期文书为定位起点；官修记录、奏报或外交文本服务特定机构立场，具体日期、规模、地方执行与对手视角须以独立材料复核。
\hypertarget{social:EQ-1861-XIANFENG-DEATH}{}\paragraph{咸丰十一年（1861年）}咸丰帝去世，同治帝即位，顾命大臣与两宫太后权力冲突爆发。从地方治理、身份、家庭与物质生活的角度，这一节点要求追问它在何处改变了资源、道路、组织或日常风险。核心来源为《清文宗实录》咸丰十一年相关条；宫中朱批奏折、内阁大库题本与《清史稿》相关志传。它首先留下的是官府分类、任命、处置或资源调拨的痕迹；要判断基层是否照办，仍须寻找县档、账簿、契约或后续案件。账本所示的研究方向是宫廷政治、制度运作与中央—地方信息链研究。可核的条件包括战时政局与幼主继承形成权力真空。应分层观察的后果是辛酉政变重塑清廷中枢决策结构。证据界限：以咸丰十一年（1861年）的同期文书为定位起点；官修记录、奏报或外交文本服务特定机构立场，具体日期、规模、地方执行与对手视角须以独立材料复核。
\hypertarget{social:EQ-1861-XINYOU-COUP}{}\paragraph{咸丰十一年（1861年）}慈禧、慈安与恭亲王发动辛酉政变，肃顺等顾命大臣失势。从地方治理、身份、家庭与物质生活的角度，这一节点要求追问它在何处改变了资源、道路、组织或日常风险。核心来源为《清文宗实录》咸丰十一年相关条；宫中朱批奏折、内阁大库题本与《清史稿》相关志传。它首先留下的是官府分类、任命、处置或资源调拨的痕迹；要判断基层是否照办，仍须寻找县档、账簿、契约或后续案件。账本所示的研究方向是宫廷政治、制度运作与中央—地方信息链研究。可核的条件包括幼主摄政安排与京师政治网络冲突。应分层观察的后果是两宫垂帘与恭亲王参与中枢成为同治初年格局。证据界限：以咸丰十一年（1861年）的同期文书为定位起点；官修记录、奏报或外交文本服务特定机构立场，具体日期、规模、地方执行与对手视角须以独立材料复核。
\hypertarget{social:EQ-1861-ZONGLI-YAMEN}{}\paragraph{咸丰十一年（1861年）}总理各国事务衙门设立，清廷开始形成常设近代外交处理机构。从地方治理、身份、家庭与物质生活的角度，这一节点要求追问它在何处改变了资源、道路、组织或日常风险。核心来源为《清文宗实录》咸丰十一年相关条；《筹办夷务始末》相关年卷与中外照会、条约文本。它首先留下的是官府分类、任命、处置或资源调拨的痕迹；要判断基层是否照办，仍须寻找县档、账簿、契约或后续案件。账本所示的研究方向是条约文本、外交档案、口岸社会与国际法实践研究。可核的条件包括条约体系和驻京使节要求持续对外交涉能力。应分层观察的后果是机构职责不断调整，不能视为完整现代外交部。证据界限：以咸丰十一年（1861年）的同期文书为定位起点；官修记录、奏报或外交文本服务特定机构立场，具体日期、规模、地方执行与对手视角须以独立材料复核。
\hypertarget{social:EQ-1862-SELF-STRENGTHENING-ORIGINS}{}\paragraph{同治元年（1862年）}洋务实践在军械、翻译、外交和地方军政中展开。从地方治理、身份、家庭与物质生活的角度，这一节点要求追问它在何处改变了资源、道路、组织或日常风险。核心来源为《清穆宗实录》同治元年相关条；宫中朱批奏折、内阁大库题本与《清史稿》相关志传。它首先留下的是官府分类、任命、处置或资源调拨的痕迹；要判断基层是否照办，仍须寻找县档、账簿、契约或后续案件。账本所示的研究方向是宫廷政治、制度运作与中央—地方信息链研究。可核的条件包括内战与列强压力促使官员寻求技术和制度调整。应分层观察的后果是不同地区的洋务项目具有不同资金、目标和绩效。证据界限：以同治元年（1862年）的同期文书为定位起点；官修记录、奏报或外交文本服务特定机构立场，具体日期、规模、地方执行与对手视角须以独立材料复核。
\hypertarget{social:EQ-1862-SHANGHAI-FOREIGN-ARMS}{}\paragraph{同治元年（1862年）}上海的军火、雇佣力量和租界网络影响清军对太平军作战。从地方治理、身份、家庭与物质生活的角度，这一节点要求追问它在何处改变了资源、道路、组织或日常风险。核心来源为《清穆宗实录》同治元年相关条；军机处档、战事方略及地方奏报。这类编年材料可以标定朝廷获知和叙述事件的时间，却往往压低妇女、儿童、佃户、流民和非文书精英的行动。账本所示的研究方向是战争动员、军需、战场暴力与地方社会研究。可核的条件包括江南战场接近条约口岸与国际商业网络。应分层观察的后果是外来军事技术和人员角色须避免夸大为决定性单因。证据界限：以同治元年（1862年）的同期文书为定位起点；官修记录、奏报或外交文本服务特定机构立场，具体日期、规模、地方执行与对手视角须以独立材料复核。
\hypertarget{social:EQ-1863-EVER-VICTORIOUS-ARMY}{}\paragraph{同治二年（1863年）}常胜军在江南战场参与清军作战，戈登接掌后扩大影响。从地方治理、身份、家庭与物质生活的角度，这一节点要求追问它在何处改变了资源、道路、组织或日常风险。核心来源为《清穆宗实录》同治二年相关条；军机处档、战事方略及地方奏报。这类编年材料可以标定朝廷获知和叙述事件的时间，却往往压低妇女、儿童、佃户、流民和非文书精英的行动。账本所示的研究方向是战争动员、军需、战场暴力与地方社会研究。可核的条件包括上海周边的财政、外商和地方军政资源可被军事化。应分层观察的后果是常胜军规模有限，其作用需放入湘淮军和地方战场比较。证据界限：以同治二年（1863年）的同期文书为定位起点；官修记录、奏报或外交文本服务特定机构立场，具体日期、规模、地方执行与对手视角须以独立材料复核。
\hypertarget{social:EQ-1863-NIAN-WAR-MOBILIZATION}{}\paragraph{同治二年（1863年）}清廷调集湘淮军和地方力量应对捻军，华北战事加剧。从地方治理、身份、家庭与物质生活的角度，这一节点要求追问它在何处改变了资源、道路、组织或日常风险。核心来源为《清穆宗实录》同治二年相关条；军机处档、战事方略及地方奏报。这类编年材料可以标定朝廷获知和叙述事件的时间，却往往压低妇女、儿童、佃户、流民和非文书精英的行动。账本所示的研究方向是战争动员、军需、战场暴力与地方社会研究。可核的条件包括太平战场与北方流动战场同时消耗军费。应分层观察的后果是军队调动改变省际财政与地方权力关系。证据界限：以同治二年（1863年）的同期文书为定位起点；官修记录、奏报或外交文本服务特定机构立场，具体日期、规模、地方执行与对手视角须以独立材料复核。
\hypertarget{social:EQ-1864-TIANJING-FALL}{}\paragraph{同治三年（1864年）}湘军攻陷天京，太平天国核心政权覆亡。从地方治理、身份、家庭与物质生活的角度，这一节点要求追问它在何处改变了资源、道路、组织或日常风险。核心来源为《清穆宗实录》同治三年相关条；军机处档、战事方略及地方奏报。这类编年材料可以标定朝廷获知和叙述事件的时间，却往往压低妇女、儿童、佃户、流民和非文书精英的行动。账本所示的研究方向是战争动员、军需、战场暴力与地方社会研究。可核的条件包括长期围困、粮食危机和内部裂解削弱天京。应分层观察的后果是战后屠杀、人口损失和城市重建须以多种材料谨慎处理。证据界限：以同治三年（1864年）的同期文书为定位起点；官修记录、奏报或外交文本服务特定机构立场，具体日期、规模、地方执行与对手视角须以独立材料复核。
\hypertarget{social:EQ-1865-DUNGAN-REVOLT}{}\paragraph{同治四年（1865年）}西北回民起事在陕西、甘肃等地扩展，形成长期战争。从地方治理、身份、家庭与物质生活的角度，这一节点要求追问它在何处改变了资源、道路、组织或日常风险。核心来源为《清穆宗实录》同治四年相关条；军机处档、战事方略及地方奏报。这类编年材料可以标定朝廷获知和叙述事件的时间，却往往压低妇女、儿童、佃户、流民和非文书精英的行动。账本所示的研究方向是战争动员、军需、战场暴力与地方社会研究。可核的条件包括地方冲突、军政失序和族群分类相互激化。应分层观察的后果是动机、伤亡和社区边界须用地方与多语材料审慎叙述。证据界限：以同治四年（1865年）的同期文书为定位起点；官修记录、奏报或外交文本服务特定机构立场，具体日期、规模、地方执行与对手视角须以独立材料复核。
\hypertarget{social:EQ-1865-NIAN-LEADER-DEATH}{}\paragraph{同治四年（1865年）}张乐行等捻军领导人先后失势或被杀，捻军仍持续活动。从地方治理、身份、家庭与物质生活的角度，这一节点要求追问它在何处改变了资源、道路、组织或日常风险。核心来源为《清穆宗实录》同治四年相关条；军机处档、战事方略及地方奏报。这类编年材料可以标定朝廷获知和叙述事件的时间，却往往压低妇女、儿童、佃户、流民和非文书精英的行动。账本所示的研究方向是战争动员、军需、战场暴力与地方社会研究。可核的条件包括清军加强围堵与地方武装协作。应分层观察的后果是领导层变化不等于流动作战网络立即瓦解。证据界限：以同治四年（1865年）的同期文书为定位起点；官修记录、奏报或外交文本服务特定机构立场，具体日期、规模、地方执行与对手视角须以独立材料复核。
\hypertarget{social:EQ-1866-FUJIAN-SHIPYARD}{}\paragraph{同治五年（1866年）}福州船政局创办，成为洋务造船、教育与技术翻译的重要项目。从地方治理、身份、家庭与物质生活的角度，这一节点要求追问它在何处改变了资源、道路、组织或日常风险。核心来源为《清穆宗实录》同治五年相关条；户部奏销、盐政、河工或海关档案。它首先留下的是官府分类、任命、处置或资源调拨的痕迹；要判断基层是否照办，仍须寻找县档、账簿、契约或后续案件。账本所示的研究方向是财政账册、市场网络、灾害环境与区域经济研究。可核的条件包括海防压力和技术学习需求增长。应分层观察的后果是厂局经费、产能和军事效果需用账册与具体船只记录评估。证据界限：以同治五年（1866年）的同期文书为定位起点；官修记录、奏报或外交文本服务特定机构立场，具体日期、规模、地方执行与对手视角须以独立材料复核。
\hypertarget{social:EQ-1866-NINGBO-SHANGHAI-INDUSTRY}{}\paragraph{同治五年（1866年）}江南制造局等军工项目在战后江南环境中继续发展。从地方治理、身份、家庭与物质生活的角度，这一节点要求追问它在何处改变了资源、道路、组织或日常风险。核心来源为《清穆宗实录》同治五年相关条；户部奏销、盐政、河工或海关档案。它首先留下的是官府分类、任命、处置或资源调拨的痕迹；要判断基层是否照办，仍须寻找县档、账簿、契约或后续案件。账本所示的研究方向是财政账册、市场网络、灾害环境与区域经济研究。可核的条件包括湘淮军和条约口岸带来技术、资金与人才条件。应分层观察的后果是“自强”口号不能替代对生产规模和依赖进口的检验。证据界限：以同治五年（1866年）的同期文书为定位起点；官修记录、奏报或外交文本服务特定机构立场，具体日期、规模、地方执行与对手视角须以独立材料复核。
\hypertarget{social:EQ-1867-NIAN-CAVALRY-CAMPAIGN}{}\paragraph{同治六年（1867年）}清军以河防、堡垒和机动军队围堵捻军骑兵。从地方治理、身份、家庭与物质生活的角度，这一节点要求追问它在何处改变了资源、道路、组织或日常风险。核心来源为《清穆宗实录》同治六年相关条；军机处档、战事方略及地方奏报。这类编年材料可以标定朝廷获知和叙述事件的时间，却往往压低妇女、儿童、佃户、流民和非文书精英的行动。账本所示的研究方向是战争动员、军需、战场暴力与地方社会研究。可核的条件包括捻军流动作战突破传统省界防御。应分层观察的后果是军事策略的效果应与地方破坏和民众负担并看。证据界限：以同治六年（1867年）的同期文书为定位起点；官修记录、奏报或外交文本服务特定机构立场，具体日期、规模、地方执行与对手视角须以独立材料复核。
\hypertarget{social:EQ-1868-NIAN-SUPPRESSION}{}\paragraph{同治七年（1868年）}东捻军等主力被击败，捻军大规模战争告终。从地方治理、身份、家庭与物质生活的角度，这一节点要求追问它在何处改变了资源、道路、组织或日常风险。核心来源为《清穆宗实录》同治七年相关条；军机处档、战事方略及地方奏报。这类编年材料可以标定朝廷获知和叙述事件的时间，却往往压低妇女、儿童、佃户、流民和非文书精英的行动。账本所示的研究方向是战争动员、军需、战场暴力与地方社会研究。可核的条件包括淮军等部队的围堵和战场分割发挥作用。应分层观察的后果是战后地方武装和财政动员遗产继续存在。证据界限：以同治七年（1868年）的同期文书为定位起点；官修记录、奏报或外交文本服务特定机构立场，具体日期、规模、地方执行与对手视角须以独立材料复核。
\hypertarget{social:EQ-1869-YANGZHOU-SALT-CRISIS}{}\paragraph{同治八年（1869年）}两淮盐政改革、商欠与财政重整成为中兴时期的重要议题。从地方治理、身份、家庭与物质生活的角度，这一节点要求追问它在何处改变了资源、道路、组织或日常风险。核心来源为《清穆宗实录》同治八年相关条；户部奏销、盐政、河工或海关档案。它首先留下的是官府分类、任命、处置或资源调拨的痕迹；要判断基层是否照办，仍须寻找县档、账簿、契约或后续案件。账本所示的研究方向是财政账册、市场网络、灾害环境与区域经济研究。可核的条件包括战乱破坏盐业网络，厘金与新税制改变利益结构。应分层观察的后果是改革效果需比较名义制度、实收与商人账目。证据界限：以同治八年（1869年）的同期文书为定位起点；官修记录、奏报或外交文本服务特定机构立场，具体日期、规模、地方执行与对手视角须以独立材料复核。
\subsection{1870—1879：地方材料所能看见的尺度}
以下事件按账本时间排列；它们不是同质的社会样本，而是提示应到何种地点、机构和材料中继续追问的入口。
\hypertarget{social:EQ-1870-TIANJIN-MASSACRE}{}\paragraph{同治九年（1870年）}天津教案发生，法国领事、修女和中国民众死伤，引发外交危机。从地方治理、身份、家庭与物质生活的角度，这一节点要求追问它在何处改变了资源、道路、组织或日常风险。核心来源为《清穆宗实录》同治九年相关条；《筹办夷务始末》相关年卷与中外照会、条约文本。它首先留下的是官府分类、任命、处置或资源调拨的痕迹；要判断基层是否照办，仍须寻找县档、账簿、契约或后续案件。账本所示的研究方向是条约文本、外交档案、口岸社会与国际法实践研究。可核的条件包括教会活动、地方谣言和条约保护关系高度紧张。应分层观察的后果是责任叙述在中法档案中差异显著，须避免单一归因。证据界限：以同治九年（1870年）的同期文书为定位起点；官修记录、奏报或外交文本服务特定机构立场，具体日期、规模、地方执行与对手视角须以独立材料复核。
\hypertarget{social:EQ-1870-ZHILI-FAMINE-RELIEF}{}\paragraph{同治九年（1870年）}华北灾荒和赈务问题促使官员、士绅和国际救济网络介入。从地方治理、身份、家庭与物质生活的角度，这一节点要求追问它在何处改变了资源、道路、组织或日常风险。核心来源为《清穆宗实录》同治九年相关条；地方志、督抚奏折及地方档案。这类编年材料可以标定朝廷获知和叙述事件的时间，却往往压低妇女、儿童、佃户、流民和非文书精英的行动。账本所示的研究方向是地方档案、迁徙、宗教、族群与社会动员研究。可核的条件包括旱灾、粮价、交通和地方储备共同影响灾情。应分层观察的后果是慈善报告与官府奏报的覆盖范围和数字需核对。证据界限：以同治九年（1870年）的同期文书为定位起点；官修记录、奏报或外交文本服务特定机构立场，具体日期、规模、地方执行与对手视角须以独立材料复核。
\hypertarget{social:EQ-1871-MUDAN-INCIDENT}{}\paragraph{同治十年（1871年）}琉球船民在台湾牡丹社一带遇害，事件成为日本对台政策借口。从地方治理、身份、家庭与物质生活的角度，这一节点要求追问它在何处改变了资源、道路、组织或日常风险。核心来源为《清穆宗实录》同治十年相关条；《筹办夷务始末》相关年卷与中外照会、条约文本。它首先留下的是官府分类、任命、处置或资源调拨的痕迹；要判断基层是否照办，仍须寻找县档、账簿、契约或后续案件。账本所示的研究方向是条约文本、外交档案、口岸社会与国际法实践研究。可核的条件包括海难、原住民地域与清朝地方治理边界交错。应分层观察的后果是事件叙述须区分原住民行动、清方管辖与日方外交利用。证据界限：以同治十年（1871年）的同期文书为定位起点；官修记录、奏报或外交文本服务特定机构立场，具体日期、规模、地方执行与对手视角须以独立材料复核。
\hypertarget{social:EQ-1871-QING-STUDENTS-ABROAD}{}\paragraph{同治十年（1871年）}清廷派遣幼童赴美留学的计划启动，洋务教育进入海外培训阶段。从地方治理、身份、家庭与物质生活的角度，这一节点要求追问它在何处改变了资源、道路、组织或日常风险。核心来源为《清穆宗实录》同治十年相关条；内阁档、文集或报刊相关记录。这类编年材料可以标定朝廷获知和叙述事件的时间，却往往压低妇女、儿童、佃户、流民和非文书精英的行动。账本所示的研究方向是知识生产、教育、宗教与公共传播研究。可核的条件包括技术、语言与外交人才需求上升。应分层观察的后果是选派和留学经历不代表全国教育制度已转型。证据界限：以同治十年（1871年）的同期文书为定位起点；官修记录、奏报或外交文本服务特定机构立场，具体日期、规模、地方执行与对手视角须以独立材料复核。
\hypertarget{social:EQ-1872-CHINA-MERCHANTS-STEAM-NAVIGATION}{}\paragraph{同治十一年（1872年）}轮船招商局创办，试图以官督商办方式发展近代航运。从地方治理、身份、家庭与物质生活的角度，这一节点要求追问它在何处改变了资源、道路、组织或日常风险。核心来源为《清穆宗实录》同治十一年相关条；户部奏销、盐政、河工或海关档案。它首先留下的是官府分类、任命、处置或资源调拨的痕迹；要判断基层是否照办，仍须寻找县档、账簿、契约或后续案件。账本所示的研究方向是财政账册、市场网络、灾害环境与区域经济研究。可核的条件包括外轮竞争和沿海、长江运输需求推动新企业。应分层观察的后果是企业性质、资本来源和经营绩效需依账册分析。证据界限：以同治十一年（1872年）的同期文书为定位起点；官修记录、奏报或外交文本服务特定机构立场，具体日期、规模、地方执行与对手视角须以独立材料复核。
\hypertarget{social:EQ-1873-TONGZHI-EMPEROR-REIGN}{}\paragraph{同治十二年（1873年）}同治帝亲政，垂帘体制与皇帝个人政治角色发生调整。从地方治理、身份、家庭与物质生活的角度，这一节点要求追问它在何处改变了资源、道路、组织或日常风险。核心来源为《清穆宗实录》同治十二年相关条；宫中朱批奏折、内阁大库题本与《清史稿》相关志传。它首先留下的是官府分类、任命、处置或资源调拨的痕迹；要判断基层是否照办，仍须寻找县档、账簿、契约或后续案件。账本所示的研究方向是宫廷政治、制度运作与中央—地方信息链研究。可核的条件包括幼帝成年改变摄政安排。应分层观察的后果是实际决策仍受太后、军机和重臣网络制约。证据界限：以同治十二年（1873年）的同期文书为定位起点；官修记录、奏报或外交文本服务特定机构立场，具体日期、规模、地方执行与对手视角须以独立材料复核。
\hypertarget{social:EQ-1873-JAPAN-RYUKYU-DISPUTE}{}\paragraph{同治十二年（1873年）}日本围绕琉球与台湾事务加强对清交涉。从地方治理、身份、家庭与物质生活的角度，这一节点要求追问它在何处改变了资源、道路、组织或日常风险。核心来源为《清穆宗实录》同治十二年相关条；《筹办夷务始末》相关年卷与中外照会、条约文本。它首先留下的是官府分类、任命、处置或资源调拨的痕迹；要判断基层是否照办，仍须寻找县档、账簿、契约或后续案件。账本所示的研究方向是条约文本、外交档案、口岸社会与国际法实践研究。可核的条件包括日本明治国家扩张与东亚朝贡关系重组。应分层观察的后果是册封、属国和主权词汇应置于各方文本语境中。证据界限：以同治十二年（1873年）的同期文书为定位起点；官修记录、奏报或外交文本服务特定机构立场，具体日期、规模、地方执行与对手视角须以独立材料复核。
\hypertarget{social:EQ-1874-JAPANESE-TAIWAN-EXPEDITION}{}\paragraph{同治十三年（1874年）}日本以牡丹社事件为由出兵台湾南部，清廷被迫应对。从地方治理、身份、家庭与物质生活的角度，这一节点要求追问它在何处改变了资源、道路、组织或日常风险。核心来源为《清穆宗实录》同治十三年相关条；军机处档、战事方略及地方奏报。这类编年材料可以标定朝廷获知和叙述事件的时间，却往往压低妇女、儿童、佃户、流民和非文书精英的行动。账本所示的研究方向是战争动员、军需、战场暴力与地方社会研究。可核的条件包括清廷对台湾东部和原住民地区的治理能力有限。应分层观察的后果是战后赔款与外交安排刺激清廷加强台湾经营。证据界限：以同治十三年（1874年）的同期文书为定位起点；官修记录、奏报或外交文本服务特定机构立场，具体日期、规模、地方执行与对手视角须以独立材料复核。
\hypertarget{social:EQ-1874-TONGZHI-DEATH}{}\paragraph{同治十三年（1874年）}同治帝去世，慈禧选择光绪帝继位，再次形成幼主政治。从地方治理、身份、家庭与物质生活的角度，这一节点要求追问它在何处改变了资源、道路、组织或日常风险。核心来源为《清穆宗实录》同治十三年相关条；宫中朱批奏折、内阁大库题本与《清史稿》相关志传。它首先留下的是官府分类、任命、处置或资源调拨的痕迹；要判断基层是否照办，仍须寻找县档、账簿、契约或后续案件。账本所示的研究方向是宫廷政治、制度运作与中央—地方信息链研究。可核的条件包括同治无嗣使宗法与宫廷权力问题交织。应分层观察的后果是慈禧重新垂帘，晚清中枢进入新阶段。证据界限：以同治十三年（1874年）的同期文书为定位起点；官修记录、奏报或外交文本服务特定机构立场，具体日期、规模、地方执行与对手视角须以独立材料复核。
